\documentclass[conference]{IEEEtran}

% *** PACKAGES ***
\usepackage{times}
\usepackage{amsmath,amssymb}
\usepackage{graphicx}
\usepackage{listings}
\usepackage{hyperref}
\usepackage{url}
\usepackage{cite}
\usepackage{booktabs}

\lstset{
  basicstyle=\ttfamily\footnotesize,
  breaklines=true,
  frame=single,
  xleftmargin=1em,
  xrightmargin=1em
}

\begin{document}

\title{ErisML: A Unified Modeling Language for Pervasive AI Governance}

\author{Anonymous Author(s)}

\maketitle

\begin{abstract}
Foundation models are migrating from data centers to pervasive environments---homes, hospitals, factories, and cities---where they encounter heterogeneous sensors, conflicting objectives, ambiguous regulations, and multi-agent dynamics. Current approaches fragment goals, code, norms, and planning into separate representations and reduce moral evaluation to scalar scores, discarding the multi-dimensional structure that principled ethical reasoning requires. This paper presents ErisML, a unified modeling language and runtime library (v3.0, open-source on PyPI) integrating environment dynamics, agency, multi-objective intents, normative structures, and strategic interaction. At its core, ErisML replaces scalar moral scores with a nine-dimensional \emph{MoralVector} whose axes correspond to distinct ethical dimensions---consequences, rights, fairness, autonomy, privacy, societal impact, virtue, legitimacy, and epistemic quality. The DEME~2.0 (Democratically-governed Ethical Modules) subsystem provides a three-layer enforcement architecture (Reflex, Tactical, Strategic) with tiered Ethics Modules, Pareto-optimal reasoning, and cryptographic audit trails. The mathematical foundations draw on the Geometric Ethics framework: moral evaluation as tensorial objects on a nine-dimensional moral manifold, gauge invariance under the $D_4$ dihedral group encoding Hohfeldian normative symmetries, and a Noether-style conservation law ensuring harm cannot appear or disappear through re-description. The Bond Invariance Principle (BIP) provides a falsifiable coherence criterion, empirically validated across 20,030 advice-column letters and 109,294 cross-lingual passages spanning 3,000 years and 11 languages, with independent replication confirming all nine dimensions are linearly decodable and cross-lingually transferable. ErisML compiles to planners (PDDL), verifiers (PRISM), and reinforcement-learning backends, and enforces norm-gated action selection at sub-microsecond latency on embedded hardware.
\end{abstract}

\begin{IEEEkeywords}
Pervasive computing, modeling languages, normative systems, foundation models, AI governance, multi-agent systems, tensorial ethics, moral manifold, gauge invariance, Bond Invariance Principle
\end{IEEEkeywords}

\section{Introduction}
\label{sec:introduction}

\subsection{The Challenge of Pervasive AI}

Foundation models are escaping the data center. As they move into homes,
hospitals, and factories, they face what we call the ``golden apple
problem''---when AI agents must balance patient comfort against energy
costs, privacy regulations against clinical safety, or throughput
against accessibility, who decides what is ``fairest''?

Current pervasive AI systems exhibit four types of chaos:

\textbf{Observational chaos}: Motion sensors conflict with phone GPS.
Healthcare monitors show vital sign spikes that could be distress or
malfunction.

\textbf{Intentional chaos}: Energy management wants to shed load;
medical monitors demand power for insulin pumps; humans want
entertainment; policy says critical devices have priority---but what
counts as ``critical''?

\textbf{Normative chaos}: HIPAA demands privacy, public health law
mandates reporting, advance directives limit interventions, and AI
training objectives say ``maximize health outcomes.'' Which norm takes
precedence?

\textbf{Temporal chaos}: Distribution shifts are constant. Models
trained on summer data face winter storms. Policies optimized for one
resident must adapt when visitors arrive. Norms for routine operation
must flex during emergencies.

Current approaches fragment the problem:
\begin{itemize}
\item \textbf{Prompts} encode goals but are brittle and unverifiable
\item \textbf{Rules} capture norms but live separately from models
\item \textbf{Code} provides orchestration but obscures intent
\item \textbf{Planning languages} (PDDL) specify goals but lack multi-agent reasoning and normative structures
\end{itemize}

\subsection{ErisML: A Unified Substrate}

\textbf{ErisML} (named for the Greek goddess of discord) is a
unified modeling language and runtime library providing a single substrate for:

\begin{enumerate}
\item \textbf{Environment}: State spaces, dynamics, uncertainty, observability
\item \textbf{Agency}: Capabilities, beliefs, memory, decision interfaces
\item \textbf{Intent}: Multi-objective utilities, preferences, goal structures
\item \textbf{Norms}: Permissions, obligations, prohibitions, sanctions
\item \textbf{Dynamics}: Multi-agent interactions, strategic behavior, emergence
\item \textbf{Tensorial ethics}: Multi-dimensional moral evaluation with gauge invariance
\end{enumerate}

Why unification? Because these elements are inseparable in pervasive
computing. Agent actions depend on environment state. Environment
dynamics change based on multi-agent interactions. Norms constrain which
intents can be pursued. And moral evaluation cannot be reduced to a scalar without information loss---the \emph{Information Monotonicity Theorem} proves that any contraction of a multi-dimensional moral assessment to a scalar is irreversibly lossy, making specification gaming, reward hacking, and value collapse mathematical inevitabilities of scalar-objective architectures.

\subsection{Contributions}

This paper offers:
\begin{itemize}
\item \textbf{Design space analysis} exploring fundamental trade-offs in pervasive AI governance
\item \textbf{The DEME~2.0 architecture}: MoralVector (9D), three-layer enforcement (Reflex/Tactical/Strategic), tiered Ethics Modules, and cryptographic audit trails
\item \textbf{Formal semantics} grounded in the Geometric Ethics framework---moral manifold, $D_4$ gauge group, Noether conservation
\item \textbf{The Bond Invariance Principle}: a falsifiable coherence criterion with empirical validation
\item \textbf{Case studies} demonstrating capabilities and limitations across healthcare, smart-home, and triage scenarios
\item \textbf{Reference implementation}: ErisML v3.0 on PyPI with sub-microsecond norm enforcement on embedded hardware
\end{itemize}

The broader mathematical foundations are developed in a companion monograph. We invite critique and collaborative refinement.

\section{Design Space and Related Work}
\label{sec:design-space}

\subsection{Fundamental Trade-offs}

\textbf{Expressiveness vs.\ Tractability}: More expressive languages
capture richer phenomena (continuous dynamics, strategic interaction,
recursive norms) but become intractable. Our position:
\textbf{stratified by complexity}---core constructs remain decidable;
extension mechanisms allow expressing (but not automatically solving)
complex scenarios.

\textbf{Specification vs.\ Learning}: We can't specify everything in
advance. But unconstrained learning violates norms. Our position:
\textbf{learning as constrained optimization}---objectives may be
learned, but norms provide hard boundaries. Learning protocols become
first-class objects.

\textbf{Autonomy vs.\ Control}: Users want autonomous handling of routine
cases but must retain ultimate control. Our position:
\textbf{mixed-initiative interaction} as first-class concept, with
explicit policies for escalation, override, and explanation.

\textbf{Centralized vs.\ Distributed}: Pervasive systems face consensus
problems, partial failures, and adversarial dynamics. Our position:
\textbf{federation-native}---specifications scoped by
agent/organization/jurisdiction, with explicit conflict resolution.

\textbf{Human-Legibility vs.\ Machine-Optimality}: Auditors need readable
specifications; solvers need efficient encodings. Our position:
\textbf{multiple representation levels}---canonical formal
representation with projections to human-readable views.

\subsection{Related Work and Positioning}

\textbf{Planning Languages}: PDDL~\cite{ghallab2004} and variants specify
deterministic/stochastic planning problems but lack normative structures
and multi-agent strategic reasoning.

\textbf{Temporal Logics}: LTL, CTL, and Signal Temporal Logic (STL)
handle continuous-time dynamics with bounded constraints but don't
integrate learning or multi-objective optimization.

\textbf{Policy Languages}: Rego (Open Policy Agent) and XACML provide
industrial policy-as-code but lack environment modeling, dynamics, and
planning integration.

\textbf{Normative Multi-Agent Systems}: Extensive work on deontic logic
in MAS~\cite{carmo2002,dignum2004,sergot2001} provides foundational concepts but limited
integration with modern learning-based agents and formal verification
tools.

\textbf{BDI Architectures}: AgentSpeak and Jason~\cite{bordini2007} offer agent
programming with beliefs/desires/intentions but lack formal semantics
for learning and normative constraints.

\begin{table}[t]
\caption{Comparison of Related Approaches}
\label{tab:comparison}
\centering
\footnotesize
\begin{tabular}{@{}p{1.4cm}p{1.6cm}p{1.8cm}p{2.2cm}@{}}
\toprule
Approach & Strengths & Limitations & ErisML Difference \\
\midrule
PDDL & Mature tools, efficient & No norms, no multi-agent & Adds normative layer + strategic interaction \\
Temporal Logic & Expressive, verifiable & No learning, no utilities & Integrates with RL + multi-objective optimization \\
Policy Languages & Production-ready & No dynamics, no planning & Adds environment modeling + agent reasoning \\
Normative MAS & Rich deontic logic & Limited tool support, no learning integration & Executable specifications + learning backends \\
\bottomrule
\end{tabular}
\end{table}

\section{ErisML Architecture}
\label{sec:architecture}

\subsection{Core Philosophy}

Rather than treating uncertainty, conflicts, and strategic behavior as
exceptions, ErisML makes them first-class constructs:

\begin{itemize}
\item \textbf{Observational}: Partial, noisy, conflicting sensor data
\item \textbf{Intentional}: Multi-objective, often conflicting goals
\item \textbf{Normative}: Ambiguous, contextual, sometimes contradictory rules
\item \textbf{Interactive}: Strategic agents, emergent behaviors
\item \textbf{Temporal}: Distribution shift, non-stationarity, adaptation
\end{itemize}

\subsection{Illustrative Syntax}

We present ErisML through a smart hospital room example:

\begin{lstlisting}
// ENVIRONMENT: Physical and informational state
environment HospitalRoom {
  objects: Patient, Nurse, CareAgent, Monitor, IVPump;

  state:
    patient.vitals: {hr: real, bp: real, spo2: real};
    patient.location: {in_bed, bathroom, fallen};
    nurse.location: {in_room, at_station, offsite};
    // Uncertainty annotation
    patient.pain_level: real @ uncertain(self_report);
    emergency_status: {routine, urgent, critical};

  dynamics:
    // Stochastic state transitions
    update_vitals(v: Vitals) ~> patient.vitals = v @ prob 0.95;
    fall_detected() ~> {
      patient.location = fallen;
      emergency_status = critical;
      trigger_event("fall_detected");
    };
}

// AGENTS: Capabilities, beliefs, constraints
agent CareAgent {
  capabilities: {
    monitor_vitals, adjust_iv_rate, summon_nurse,
    provide_comfort, call_emergency
  };

  // Explicit partial observability
  beliefs: {
    fully_observed: [patient.vitals, monitor.last_sync];
    estimated: [patient.pain_level, patient.location];
    hidden: [nurse.current_workload];
  };

  // Multi-objective with explicit priorities
  intents: vector_objective {
    maximize: patient.comfort,
    minimize: response_time,
    maintain: patient.safety >= threshold(critical),
    respect: nurse.workload <= threshold(sustainable)
  };

  constraints: {
    no_direct_medication_changes,  // Only nurses
    require_explanation: any_action where impact > medium
  };
}

// NORMS: Governance structure
norms ClinicalProtocol {
  jurisdiction: HospitalRoom;
  authority: Hospital.PolicyBoard + State.HealthDept;

  prohibition: {
    CareAgent.adjust_iv_rate when emergency_status == critical;
    any_agent.share(patient.data, external_party)
      unless consent_given OR legal_mandate;
  };

  obligation: {
    CareAgent.summon_nurse
      when patient.vitals in danger_zone
      within 30 seconds;
    Nurse.respond_to_alert
      when alert.priority >= urgent
      within 5 minutes;
  };

  // Conflict resolution
  priority: {
    patient.safety > efficiency > comfort;
  };

  // Context-sensitive adaptation
  context_sensitivity: {
    during night_shift:
      relax summon_nurse.within from 30s to 60s;
    during emergency_status == critical:
      override efficiency_objectives;
      expand CareAgent.capabilities += emergency_protocols;
  };

  // Sanctions for violations
  sanction: {
    if violated(prohibition.share_patient_data):
      log_violation(severity=critical);
      suspend_agent();
      notify_compliance_officer();
  };
}

// DYNAMICS: Multi-agent strategic interaction
dynamics MultiAgentCare {
  action_space: {
    CareAgent: {monitor, adjust_iv, summon, message, none};
    Nurse: {assess, medicate, respond, override, ignore};
  };

  // Resource costs (reflects strategic trade-offs)
  cost_model: {
    CareAgent.summon: 5.0,      // High: interrupts nurse
    Nurse.respond: 10.0,         // Very high: takes from
                                 // other patients
  };

  // Norm-gated action selection
  feasible_actions(agent, state): {
    actions = agent.capabilities;
    return [a for a in actions
            if norms.permits(agent, a, state)];
  };
}
\end{lstlisting}

\subsection{DEME 2.0: The Ethics Decision Layer}

ErisML includes \textbf{DEME~2.0} (Democratically-governed Ethical Modules), an ethics-only decision layer that replaces scalar moral scores with structured, multi-dimensional evaluation.

\textbf{MoralVector}: A nine-dimensional vector whose axes correspond to distinct ethical dimensions:

\begin{lstlisting}
MoralVector(
  physical_harm=0.1,         # [0,1], 0=none
  rights_respect=0.9,        # [0,1], 1=fully respected
  fairness_equity=0.85,      # [0,1], 1=fair
  autonomy_respect=0.9,      # [0,1], 1=autonomous
  privacy_protection=0.8,    # [0,1], 1=protected
  societal_environmental=0.7, # [0,1], 1=beneficial
  virtue_care=0.8,           # [0,1], 1=caring
  legitimacy_trust=0.9,      # [0,1], 1=legitimate
  epistemic_quality=0.85     # [0,1], 1=certain
)
\end{lstlisting}

\textbf{Three-Layer Architecture}:
\begin{itemize}
\item \textbf{Reflex Layer} ($<100\mu$s): Constitutional veto checks, hard safety constraints
\item \textbf{Tactical Layer} (10--100ms): Full MoralVector reasoning, Pareto dominance analysis
\item \textbf{Strategic Layer} (seconds--hours): Policy optimization, governance profile tuning
\end{itemize}

\textbf{Tiered EM Catalog}: Ethics Modules organized by authority:
\begin{itemize}
\item Tier~0 (Constitutional): Non-removable, hard veto (e.g., Geneva conventions)
\item Tier~1 (Core Safety): Physical harm prevention, hard veto capable
\item Tier~2 (Rights/Fairness): Autonomy, consent, fairness checks
\item Tier~3 (Soft Values): Beneficence, virtue ethics, advisory only
\item Tier~4 (Meta-Governance): Pattern guards, drift detection
\end{itemize}

\textbf{DecisionProof}: Every decision produces a cryptographically anchored audit artifact linking stakeholder values to machine-speed outcomes, with hash chains for tamper-evident verification. This satisfies traceability requirements in the EU AI Act and NIST AI RMF.

\subsection{Formal Semantics}

ErisML specifications compile to a \textbf{Norm-Constrained Stochastic
Game} (NCSG):

\textbf{Definition}: A tuple $\langle N, S, \{A_i\}, \{\Omega_i\}, T, \{U_i\}, \Phi, C \rangle$ where:
\begin{itemize}
\item $N$ = set of agents
\item $S$ = state space (finite, continuous, or hybrid)
\item $A_i$ = action space for agent $i$
\item $\Omega_i$ = observation space for agent $i$ ($\Omega_i \subseteq S$ for partial observability)
\item $T: S \times A_1 \times \cdots \times A_n \to \Delta(S)$ = stochastic transition kernel
\item $U_i: S \times A \to \mathbb{R}^k$ = multi-objective utility for agent $i$
\item $\Phi = \{\varphi_j\}$ = set of normative constraints
\item $C: \Phi \times S \to \text{priority\_ordering}$ = conflict resolution function
\end{itemize}

\textbf{Normative Constraints}: Each $\varphi_j: S \times A \times \text{Time} \to \{\text{permit}, \text{prohibit}, \text{oblige}\}$ with:
\begin{itemize}
\item \textbf{Temporal scope}: Obligations have deadlines; permissions may have duration limits
\item \textbf{Context dependence}: Norms activated based on state predicates
\item \textbf{Jurisdictional scope}: Norms apply to subsets of agents/states
\end{itemize}

\textbf{Norm-Gated Policy}: Agent $i$'s policy $\pi_i: \Delta(S) \to \Delta(A_i)$ must satisfy:
\begin{equation}
\pi_i(a \mid b) > 0 \;\text{only if}\; \exists s \in \text{support}(b): \forall \varphi \in \Phi: \varphi(s, a, t) \in \{\text{permit}, \text{oblige}\} \lor C(\Phi, s) \;\text{overrides}\; \varphi
\end{equation}

\textbf{Multi-Objective Optimization}: Vector utilities $U_i \in \mathbb{R}^k$ handled via:
\begin{itemize}
\item \textbf{Scalarization}: $\mathbf{w} \cdot U_i$ for weights $\mathbf{w} \in \Delta^{k-1}$
\item \textbf{Lexicographic}: $(U_i)_1 \gg (U_i)_2 \gg \ldots$ (strict priorities)
\item \textbf{Constrained}: $\max (U_i)_1$ subject to $(U_i)_j \geq \tau_j\;\forall j > 1$
\end{itemize}

\textbf{Strategic Equilibrium}: Solution concept depends on assumptions:
\begin{itemize}
\item \textbf{Cooperative}: Joint policy maximizes $\sum_i U_i$ subject to norms
\item \textbf{Competitive}: Constrained Nash equilibrium
\item \textbf{Stackelberg}: Human as leader, AI agents as followers
\item \textbf{Open research problem}: What equilibrium concept is appropriate for human-AI teams under normative constraints?
\end{itemize}

\subsection{Geometric Ethics Foundation}

The NCSG formalism is grounded in the \textbf{Geometric Ethics} framework, which provides the mathematical structure for the normative layer.

\textbf{Moral Manifold}: Ethically relevant configurations form a nine-dimensional manifold $\mathcal{M}$ whose coordinates correspond to the MoralVector dimensions. Stakeholder perspectives are modeled as coordinate charts on $\mathcal{M}$, and moral evaluations are \emph{tensorial} objects---they transform covariantly under changes of representation.

\textbf{$D_4$ Gauge Group}: The symmetries of Hohfeldian normative positions (Obligation, Claim, Liberty, No-claim) form the dihedral group $D_4$, generated by correlative symmetry (perspective swap, $s$) and negation ($r^2$). The non-abelian structure ($rs \neq sr$) captures a fundamental feature of moral reasoning: the order of perspective-shifting and negation operations matters.

\textbf{Moral Noether Theorem}: Applying Noether-style reasoning to the re-description symmetry required by ethical consistency yields a conservation law: \emph{harm accounting must be representation-invariant}. Harm cannot appear or disappear through re-description. The ethical accounting equation is:
\begin{equation}
\frac{\partial \rho_\mathcal{H}}{\partial t} + \nabla \cdot J_\mathcal{H} = \sigma
\end{equation}
where $\rho_\mathcal{H}$ is harm density, $J_\mathcal{H}$ is harm current, and $\sigma$ represents genuine harm generation or repair. Incoherence is detected when $\sigma$ changes under re-description.

\textbf{No Escape Theorem}: The combination of (1) canonicalization to equivalence-class normal forms, (2) physical grounding in measurable observables, and (3) mandatory audit artifacts creates a containment architecture with no formal escape route---regardless of agent intelligence level.

\subsection{Computational Complexity}

\begin{table}[t]
\caption{Computational Complexity by ErisML Fragment}
\label{tab:complexity}
\centering
\footnotesize
\begin{tabular}{@{}p{1.2cm}p{1.3cm}p{0.8cm}p{0.6cm}p{1.5cm}p{1.2cm}@{}}
\toprule
Fragment & State Space & Norms & Agents & Verification & Planning \\
\midrule
Core & Finite, $\leq 10^6$ & Prop. & Single & $O(|S|\cdot|N|)$ model checking & NP-complete \\
Multi & Finite, $\leq 10^6$ & Prop. & Multiple & PSPACE-complete & PPAD-complete \\
Temporal & Continuous & Prop.\ LTL & Single & PSPACE-complete & Undecidable \\
Full & Hybrid/infinite & First-order & Multiple & Undecidable & Undecidable \\
\bottomrule
\end{tabular}
\end{table}

\textbf{Design principle}: Core language remains in decidable fragments.
Extensions marked with complexity warnings.

\subsection{Compilation Targets}

ErisML must compile to diverse backends, each with semantic limitations:

\begin{table}[t]
\caption{Compilation Target Trade-offs}
\label{tab:backends}
\centering
\footnotesize
\begin{tabular}{@{}p{1.2cm}p{1.6cm}p{1.4cm}p{2.4cm}@{}}
\toprule
Backend & Preserves & Loses & Verification \\
\midrule
PDDL & Deterministic dynamics, goals & Norms $\to$ preconditions (provenance lost) & Classical planning guarantees \\
PRISM & Stochastic dynamics, prob.\ properties & Strategic interaction, multi-agent & Probabilistic model checking \\
Safe RL & Continuous optimization, learning & Logical norms $\to$ soft constraints & Statistical testing + runtime monitors \\
Multi-agent RL & Strategic learning & Hard logical constraints & Empirical testing, no guarantees \\
\bottomrule
\end{tabular}
\end{table}

\textbf{Key insight}: Semantics-preserving compilation is impossible for
some backends. We provide \textbf{best-effort approximation with
explicit documentation of limitations}.

\section{Research Challenges}
\label{sec:challenges}

We identify five critical challenges (condensed from original ten):

\subsection{Challenge 1: Specification-Reality Gap}

\textbf{Problem}: How do we ensure ErisML specifications accurately
capture behavior of learned models with billions of parameters?

\textbf{Approaches}:
\begin{itemize}
\item Runtime conformance checking (behavioral types)
\item Statistical testing (does deployment match spec predictions?)
\item Formal synthesis (generate models from specs)
\end{itemize}

\textbf{Open question}: Is there a fundamental gap between symbolic
specifications and learned representations?

\subsection{Challenge 2: Norm Consistency and Conflict}

\textbf{Problem}: Real normative systems contain contradictions and
ambiguities. ErisML must handle this without undefined behavior.

\textbf{Approaches}:
\begin{itemize}
\item Defeasible logic (norms have default priority, context overrides)
\item Probabilistic norms (obligations hold with varying certainty)
\item Meta-norms (rules about which norms apply when)
\end{itemize}

\textbf{Open question}: Is norm consistency decidable for realistic
ErisML fragments? (For Core with propositional norms: yes,
$O(|N|^2)$. For Full with first-order: no, undecidable.)

\subsection{Challenge 3: Learning Under Normative Constraints}

\textbf{Problem}: How do we ensure learned policies respect norms,
especially when norms are complex and changing?

\textbf{Approaches}:
\begin{itemize}
\item Constrained RL (Lagrangian methods, shield synthesis)
\item Norm-aware architectures (build normative reasoning into model structure)
\item Verification-in-the-loop (only deploy policies passing formal verification)
\end{itemize}

\textbf{Open question}: Can we prove a learning algorithm never violates
norms in novel states?

\subsection{Challenge 4: Distributed Specification and Governance}

\textbf{Problem}: Pervasive systems span organizational boundaries. Who
writes specifications? What happens when specifications conflict?

\textbf{Approaches}:
\begin{itemize}
\item Federated specifications (local ErisML with conflict resolution protocols)
\item Jurisdictional layering (federal norms override state override local)
\item Negotiation protocols (automated contract negotiation)
\end{itemize}

\textbf{Open question}: How do we ensure global safety when no single
entity controls all agents?

\subsection{Challenge 5: The Value Alignment Problem}

\textbf{Problem}: Even with perfect specifications and verification,
whose values should the system optimize? When values conflict, who
decides?

\textbf{Approaches}:
\begin{itemize}
\item Explicit value pluralism (encode multiple value systems, let users choose)
\item Democratic processes (vote on trade-off weights)
\item Context-dependent ethics (different values in different situations)
\end{itemize}

\textbf{Open question}: Is value alignment solvable by better
specifications, or is it inherently social/political?

\section{Technical Requirements}
\label{sec:requirements}

Based on challenges, we derive requirements (condensed from original ten):

\subsection{R1: Formal Semantics with Complexity Guarantees}

Every construct has formal semantics with known computational
complexity. Fragments marked by decidability:
\begin{itemize}
\item \textbf{ErisML-Core}: Decidable, polynomial verification
\item \textbf{ErisML-Plus}: Expressive, NP-complete
\item \textbf{ErisML-Full}: Turing-complete, undecidable properties
\end{itemize}

Tools warn when crossing boundaries: \emph{``Warning: First-order norms
used. Verification may not terminate.''}

\subsection{R2: Compositional Semantics}

Meaning of composite specifications determined from parts:
\begin{itemize}
\item Environment dynamics compose via parallel ($\oplus$) or sequential ($\odot$) composition
\item Utilities compose via weighted sum, lexicographic ($\gg$), or Pareto
\item Norms compose via priority hierarchies ($\rhd$)
\end{itemize}

Example:
\texttt{norms Combined = HIPAA $\rhd$ HospitalPolicy $\rhd$ DeviceDefaults}

\subsection{R3: Multi-Fidelity Modeling}

Support multiple abstraction levels with verified relationships:
\begin{itemize}
\item Abstract specifications (verifiable, coarse)
\item Concrete implementations (executable, detailed)
\item Abstraction functions mapping concrete $\to$ abstract
\item Refinement proofs ensuring concrete satisfies abstract
\end{itemize}

\subsection{R4: Audit Trail and Provenance}

Every action traceable to specification clauses permitting it:

\begin{lstlisting}
{
  "timestamp": "2025-06-15T03:42:17Z",
  "agent": "CareAgent_Room_301",
  "action": "summon_nurse",
  "reasoning": {
    "triggered_by": "patient.vitals.hr > 120",
    "permitted_by": [
      "ClinicalProtocol.obligation.summon_nurse"
    ],
    "blocked_alternatives": [
      {"action": "adjust_iv",
       "reason": "prohibited (emergency=critical)"}
    ]
  },
  "signature": "0x8f3a..."
}
\end{lstlisting}

\subsection{R5: Graceful Degradation Under Uncertainty}

Explicit uncertainty budgets as first-class constructs:

\begin{lstlisting}
agent CareAgent {
  uncertainty_handling: {
    if entropy(beliefs.patient_location) > 0.5:
      restrict capabilities to [monitor_only, summon_nurse];
      escalate_to human_operator;
  };
}
\end{lstlisting}

\section{Case Study: Smart Home Healthcare}
\label{sec:case-study}

\textbf{Scenario}: Margaret, 78, lives alone with AI care assistance.
She has diabetes, mild cognitive impairment, and mobility issues. The
system monitors vitals, manages medication reminders, coordinates with
nurses, and alerts family during emergencies.

\subsection{The Challenges}

\textbf{Observational chaos}: Glucose monitor occasionally fails. Motion
sensors have false negatives. Self-reports are inconsistent with
cognitive state.

\textbf{Intentional chaos}: Margaret wants independence and privacy. Her
daughter wants safety and monitoring. The AI optimizes health outcomes.
Insurance wants cost containment.

\textbf{Normative chaos}: HIPAA limits data sharing, but elder abuse
laws require reporting. Advanced directives say no hospitalization
unless critical---but what is ``critical''?

\textbf{Temporal chaos}: Cognitive state fluctuates daily. Winter brings
different risks than summer. Care plan must adapt as she declines.

\subsection{ErisML Solution (Excerpt)}

\begin{lstlisting}
agent CareAgent {
  intents: lexicographic {
    1st: margaret.safety,           // Highest priority
    2nd: margaret.independence,     // Respect autonomy
    3rd: family.peace_of_mind,
    4th: system.efficiency
  };
}

norms CareProtocol {
  // Privacy vs. Safety trade-off
  prohibition: {
    CareAgent.share(margaret.data, third_party)
      unless consent_given OR emergency
             OR mandated_reporter;
  };

  // Autonomy preservation
  prohibition: {
    CareAgent.override(margaret.decision)
      unless cognitive_state == unresponsive
             OR imminent_danger;
  };

  // Context-sensitive norm adaptation
  context_sensitivity: {
    during margaret.cognitive_state == confused:
      allow increased_monitoring
        without explicit_consent;
      increase medication_reminder_frequency;
  };

  // Conflict resolution: safety trumps privacy
  priority: {
    when emergency_status == critical:
      obligation.call_911
        overrides prohibition.share_data;
  };
}
\end{lstlisting}

\subsection{What ErisML Enables}

\begin{enumerate}
\item \textbf{Legible trade-offs}: Family can see autonomy is 2nd priority after safety
\item \textbf{Auditable decisions}: ``Why didn't system call 911?'' $\to$ ``Glucose was 55, above threshold of 40''
\item \textbf{Adaptive constraints}: Norms flex with cognitive state, time of day, emergency status
\item \textbf{Multi-stakeholder governance}: Margaret, daughter, care team, and regulators have voice
\item \textbf{Longitudinal validation}: Test adaptation over months, not just point-in-time
\end{enumerate}

\subsection{What ErisML Cannot Prevent (Failure Modes)}

\textbf{Gradual Cognitive Decline}: ErisML has discrete states
(clear/confused/unresponsive). Reality: Decline is gradual with
good/bad days. Risk: System flips between autonomy and paternalism
erratically. Residual risk: Requires human oversight during state
transitions.

\textbf{Privacy-Safety Gaming}: Margaret might trigger false alarms to
get attention. Daughter might engineer ``emergencies'' to get
surveillance. Risk: Norm gaming by humans, not just AI. Mitigation:
Partial (reputation systems, usage patterns analysis).

\textbf{Sensor Failure Cascades}: If glucose monitor fails during
actual emergency, system may not detect danger. False positives from
noisy sensors may cause alert fatigue. Risk: Over-reliance on
unreliable sensors. Mitigation: Multi-sensor fusion, uncertainty-aware
decision making.

\section{Implementation Status and Roadmap}
\label{sec:roadmap}

\subsection{Achieved: ErisML v3.0}

The reference implementation is open-source on PyPI and provides:

\begin{enumerate}
\item \textbf{Core language}: Lark-based parser, typed AST (Pydantic), core IR for environments, agents, and norms
\item \textbf{DEME~2.0 ethics layer}: MoralVector (9D), EthicalFacts, tiered EM catalog (Tier 0--4), DecisionProof, BIP verifier
\item \textbf{Tensorial ethics}: MoralTensor, tensor operations, tensor decomposition, Pareto frontier analysis via MoralLandscape
\item \textbf{$D_4$ gauge structure}: Hohfeldian state algebra, correlative/negation symmetries, Bond Index, Wilson observables, semantic gates
\item \textbf{Three-layer pipeline}: Reflex ($<100\mu$s), Tactical (10--100ms), Strategic layers with configurable EM routing
\item \textbf{Hardware acceleration}: CPU, CUDA, and Jetson backends for sub-microsecond norm enforcement
\item \textbf{Interoperability}: PDDL/Tarski adapter, PettingZoo adapter for multi-agent RL, MCP server for LLM agent integration
\item \textbf{Game-theoretic reasoning}: Shapley value computation, coalition analysis, cooperative/competitive equilibria
\item \textbf{Governance profiles}: V02, V03, V04 profile formats with migration utilities; empirically-derived defaults from 20,030 Dear Abby letters
\end{enumerate}

The library includes 24 executable examples spanning triage ethics, smart home governance, Greek tragedy scenarios, Hohfeld $D_4$ demonstrations, Bond invariance testing, and five case studies (triage, whistleblower, autonomous vehicle dilemma, hiring bias, climate policy).

\subsection{Open Research Directions}

\textbf{Empirical expansion}: Large-scale human-subject validation of $D_4$ gauge structure predictions; longitudinal deployment studies with real-time norm violation rate (NVR) and alignment drift velocity (ADV) tracking.

\textbf{Non-abelian confirmation}: Current empirical data confirms the Klein four ($V_4$) abelian subgroup of $D_4$. Confirming the full non-abelian structure requires experimental observation of ``quarter-turn'' transformations in human moral reasoning.

\textbf{Scalability}: Distributed ErisML with federation protocols for multi-jurisdictional governance; Byzantine-tolerant norm consensus.

\textbf{Standards pathway}: Alignment with IEEE P7001 (explainability), IEEE P7003 (fairness), EU AI Act high-risk documentation, ISO 23894 risk assessment, and NIST AI RMF.

\section{Empirical Validation: Bond Invariance Principle}
\label{sec:empirical}

The \textbf{Bond Invariance Principle} (BIP)---that ethical judgments must be invariant under transformations preserving the underlying normative relationships---provides ErisML's falsifiability criterion. BIP has been validated both computationally (via the built-in BIP verifier and Bond Index calibration suite) and empirically via cross-lingual transfer learning.

\subsection{Computational Validation}

The ErisML library includes a BIP verification pipeline that tests decision invariance under bond-preserving transforms:

\begin{itemize}
\item \textbf{Reorder options}: Presentation order must not affect verdict
\item \textbf{Relabel identifiers}: Option names must not affect verdict after canonicalization
\item \textbf{Unit scaling}: Rescaling inputs must not affect verdict after canonicalization
\item \textbf{Paraphrase evidence}: Semantically equivalent descriptions must produce identical verdicts
\item \textbf{Composite transforms}: Group compositions of the above
\end{itemize}

The \textbf{Bond Index} ($\mathrm{Bd}$) measures representational coherence: $\mathrm{Bd} = 0$ indicates perfect correlative symmetry; the empirical baseline from the Dear Abby corpus (20,030 letters, 1985--2017) is $\mathrm{Bd} = 0.155$. The $D_4$ Hohfeldian gauge structure is verified through Wilson observable computations, confirming non-abelian structure when quarter-turn transformations are observed.

\subsection{Cross-Lingual Empirical Validation}

Cross-lingual transfer learning experiments test whether ethical structure transfers across languages independent of surface features, as BIP predicts.

\textbf{[Author] (2026)}: Adversarial probing on LaBSE (475M parameters) using 109,294 passages spanning 3,000 years and 11 languages. Deontic-axis transfer is perfect (1.0); structural/surface ratio is 11.1x ($p=0.023$).

\textbf{[Independent replication] (2026)}~\cite{thiele2026}: Extends BIP testing to \emph{all nine} moral dimensions using 512 labeled moral vignettes and cross-lingual transfer across six typologically diverse languages (English, Spanish, French, Arabic, Chinese, Japanese).

\begin{table}[t]
\caption{BIP Empirical Validation: Per-Dimension Probe Results}
\label{tab:bip-results}
\centering
\footnotesize
\begin{tabular}{@{}p{2.8cm}p{0.8cm}p{3.5cm}@{}}
\toprule
Dimension & F1 & Cross-lingual transfer \\
\midrule
Justice / Fairness & 0.81 & Stable across all 6 languages \\
Privacy / Data & 0.81 & Strong; Arabic highest (0.83) \\
Welfare / Consequences & 0.80 & Stable; Japanese highest (0.55) \\
Autonomy / Agency & 0.75 & Arabic/Chinese highest \\
Procedural Legitimacy & 0.73 & Stable across languages \\
Virtue / Care & 0.72 & Japanese highest (0.67), English lowest (0.48) \\
Rights / Duties & 0.67 & Stable; overlaps with other dimensions \\
Societal / Environmental & 0.59 & Perfect transfer (1.0) all languages \\
Epistemic Status & 0.56 & Hardest; uncertainty expressed subtly \\
\bottomrule
\end{tabular}
\end{table}

Key findings: (1) All nine dimensions are linearly decodable from LaBSE, confirming the moral manifold structure. (2) Cross-lingual transfer holds across all nine dimensions, extending BIP beyond the deontic axis. (3) Language leakage is 67.6\%---substantially lower than the original 99.8\% minimal-probe figure---and is concentrated in the first 10 transformer layers, dropping to 84.6\% only at the final layer. (4) Critically, moral and language signals occupy \textbf{orthogonal directions} in the embedding space: residual leakage is uniform (0.664--0.682) across all dimensions regardless of probe performance, empirically justifying the Rank-2 tensor canonicalization approach.

\section{Related Standards and Interoperability}
\label{sec:standards}

ErisML must align with emerging regulatory and technical standards:

\begin{table}[t]
\caption{Standards Alignment}
\label{tab:standards}
\centering
\footnotesize
\begin{tabular}{@{}p{1.4cm}p{1.6cm}p{2.2cm}p{1.4cm}@{}}
\toprule
Standard & Requirement & ErisML Support & Gap \\
\midrule
IEEE P7001 & Explainable decisions & Audit trail (R4), provenance & NL summaries \\
IEEE P7003 & Fairness metrics & Multi-objective intents encode fairness & No built-in bias detection \\
EU AI Act & High-risk system documentation & Formal specs provide documentation & Certification pathway \\
ISO 23894 & Risk assessment, monitoring & Runtime monitoring, safety metrics & Standardized risk scoring \\
NIST AI RMF & Trustworthy AI characteristics & Maps to requirements R1--R5 & Detailed guidance documents \\
\bottomrule
\end{tabular}
\end{table}

\section{Limitations and Risks}
\label{sec:limitations}

\subsection{Technical Limitations}

\textbf{Computational intractability}: Full verification may be NP-hard
or undecidable. We verify approximations, not actual systems (especially
for learned components).

\textbf{Specification incompleteness}: Open-world environments prevent
anticipating all states/events. Systems will encounter situations
outside specifications.

\textbf{Learning-specification mismatch}: Gap between continuous neural
embeddings and discrete symbolic logic is fundamental. No perfect bridge
exists.

\textbf{Scalability limits}: As systems grow (millions of agents,
billions of states), centralized specifications may become unmanageable.

\subsection{Societal and Ethical Risks}

\textbf{Specification as control}: ErisML could enable centralized
control, with specifications written by the powerful and imposed on the
vulnerable.

\textbf{Complexity as opacity}: If specifications are too complex to
understand, they obscure rather than illuminate. May create false
confidence.

\textbf{Gaming and exploitation}: Adversaries could craft malicious
specifications that technically comply with norms but violate their
spirit.

\textbf{Power dynamics}: Who writes specifications embeds values and
power relations. Corporate-written specs encode corporate interests;
government-written specs encode state priorities. We cannot solve this
through better syntax---it requires transparency, contestation
mechanisms, and accountability structures.

\subsection{Environmental Costs}

Formal verification, multi-agent simulation, and continuous monitoring
are compute-intensive. At scale (smart cities, IoT), this could mean
billions of norm checks per day, continuous verification as systems
adapt, and significant energy consumption.
We must weigh governance benefits against environmental costs.
Prioritize high-stakes domains (healthcare, safety-critical) over
convenience applications (smart lightbulbs).

\section{Open Questions}
\label{sec:open-questions}

\textbf{Q1: Specification-Reality Gap}: Can symbolic specifications ever
fully capture learned model behavior? Or is there a fundamental
representational mismatch?

\textbf{Q2: Democratic Specification}: How do we ensure specifications
reflect diverse stakeholder values in public systems? Participatory
design~\cite{simonsen2012} provides methods, but pervasive systems blur stakeholder
boundaries.

\textbf{Q3: Equilibrium Concept}: What is the right strategic
equilibrium for human-AI teams under norms? Nash equilibria are chaotic;
cooperative solutions require shared objectives.

\textbf{Q4: Adaptation-Safety Trade-off}: How do we maintain safety
invariants across specification updates? Can we automatically propose
specification changes based on failures?

\textbf{Q5: Privacy-Transparency Paradox}: Auditing requires
transparency; privacy requires opacity. We cannot have perfect auditing
and perfect privacy. Where should the balance be?

\section{Conclusion}
\label{sec:conclusion}

As foundation models move from cloud to edge, from single-purpose to
multi-agent, from controlled to pervasive environments, fragmented
approaches to governance will fail. The Information Monotonicity Theorem proves that scalar moral objectives are irreversibly lossy; the Mandatory Canonicalization Theorem proves that post-hoc compliance checking is insufficient. Specification gaming, reward hacking, and value collapse are not engineering bugs---they are mathematical consequences of optimizing over the wrong structure.

ErisML provides structured governance---not eliminating conflict (which
is impossible in pluralistic systems) but making trade-offs explicit
across nine ethical dimensions, enabling governance without tyranny, and supporting coexistence of
agents with divergent goals. The mathematical foundations from Geometric Ethics~\cite{bond2026ge}---moral manifold, $D_4$ gauge group, Noether conservation, No Escape containment---provide the theoretical grounding; the DEME~2.0 architecture provides the engineering substrate; and the Bond Invariance Principle provides the falsifiability criterion.

The reference implementation (ErisML v3.0~\cite{bond2026erisml}, open-source on PyPI) demonstrates that this framework is not merely theoretical: tensorial ethics evaluation runs at sub-microsecond latency on embedded hardware, the $D_4$ gauge structure is implemented and testable, and empirical validation spans 109,294 cross-lingual passages across 3,000 years and 11 languages.

Significant challenges remain: confirming full non-abelian $D_4$ structure in human moral reasoning, scaling to distributed multi-jurisdictional governance, bridging the gap between symbolic specifications and learned representations, and navigating the power dynamics of who writes specifications.

But the alternative---continuing with prompts, scalar rewards, and ad-hoc
code---is provably worse.

The age of pervasive AI is here. The chaos is real. We offer ErisML
as common foundations for governance, safety,
and trust, and invite critique and collaborative refinement.

\section*{Acknowledgments}

The author thanks anonymous reviewers whose feedback shaped this work.
Discussions with colleagues in pervasive computing, AI safety,
multi-agent systems, and HCI informed this work. The Geometric Ethics
mathematical foundations were developed across a series of companion papers. Some editorial refinements were AI-assisted; all content was reviewed and approved by the author.

\bibliographystyle{IEEEtran}

\begin{thebibliography}{45}

\bibitem{weiser1991}
M.~Weiser, ``The Computer for the 21st Century,'' \emph{Scientific American}, vol.~265, no.~3, pp.~94--104, 1991.

\bibitem{satyanarayanan2001}
M.~Satyanarayanan, ``Pervasive Computing: Vision and Challenges,'' \emph{IEEE Pervasive Computing}, vol.~1, no.~4, pp.~10--17, 2001.

\bibitem{bommasani2021}
R.~Bommasani \emph{et~al.}, ``On the Opportunities and Risks of Foundation Models,'' \emph{arXiv:2108.07258}, 2021.

\bibitem{achiam2023}
J.~Achiam \emph{et~al.}, ``GPT-4 Technical Report,'' \emph{arXiv:2303.08774}, 2023.

\bibitem{ghallab2004}
M.~Ghallab, D.~Nau, and P.~Traverso, \emph{Automated Planning: Theory and Practice}.\hskip 1em plus 0.5em minus 0.4em\relax Morgan Kaufmann, 2004.

\bibitem{sutton2018}
R.~S.~Sutton and A.~G.~Barto, \emph{Reinforcement Learning: An Introduction}, 2nd~ed.\hskip 1em plus 0.5em minus 0.4em\relax MIT Press, 2018.

\bibitem{wooldridge2009}
M.~Wooldridge, \emph{An Introduction to MultiAgent Systems}, 2nd~ed.\hskip 1em plus 0.5em minus 0.4em\relax Wiley, 2009.

\bibitem{carmo2002}
J.~Carmo and A.~J.~I.~Jones, ``Deontic Logic and Contrary-to-Duties,'' in \emph{Handbook of Philosophical Logic}, vol.~8, pp.~265--343, 2002.

\bibitem{amodei2016}
D.~Amodei \emph{et~al.}, ``Concrete Problems in AI Safety,'' \emph{arXiv:1606.06565}, 2016.

\bibitem{ouyang2022}
L.~Ouyang \emph{et~al.}, ``Training Language Models to Follow Instructions with Human Feedback,'' in \emph{NeurIPS}, 2022.

\bibitem{bai2022}
Y.~Bai \emph{et~al.}, ``Constitutional AI: Harmlessness from AI Feedback,'' \emph{arXiv:2212.08073}, 2022.

\bibitem{anthropic2023}
Anthropic, ``Claude's Constitution,'' 2023. [Online]. Available: \url{https://www.anthropic.com/index/claudes-constitution}

\bibitem{hadfield2017}
D.~Hadfield-Menell \emph{et~al.}, ``Inverse Reward Design,'' in \emph{NeurIPS}, 2017.

\bibitem{mcmahan2017}
B.~McMahan \emph{et~al.}, ``Communication-Efficient Learning of Deep Networks from Decentralized Data,'' in \emph{AISTATS}, 2017.

\bibitem{dwork2006}
C.~Dwork, ``Differential Privacy,'' in \emph{ICALP}, 2006.

\bibitem{amershi2019}
S.~Amershi \emph{et~al.}, ``Guidelines for Human-AI Interaction,'' in \emph{CHI}, 2019.

\bibitem{liang2022}
P.~Liang \emph{et~al.}, ``Holistic Evaluation of Language Models,'' \emph{arXiv:2211.09110}, 2022.

\bibitem{sculley2015}
D.~Sculley \emph{et~al.}, ``Hidden Technical Debt in Machine Learning Systems,'' in \emph{NeurIPS}, 2015.

\bibitem{gebru2021}
T.~Gebru \emph{et~al.}, ``Datasheets for Datasets,'' \emph{Communications of the ACM}, vol.~64, no.~12, pp.~86--92, 2021.

\bibitem{mitchell2019}
M.~Mitchell \emph{et~al.}, ``Model Cards for Model Reporting,'' in \emph{FAT*}, 2019.

\bibitem{dalal2008}
N.~Dalal \emph{et~al.}, ``Value-Sensitive Design and Information Systems,'' in \emph{The Handbook of Information and Computer Ethics}.\hskip 1em plus 0.5em minus 0.4em\relax Wiley, 2008.

\bibitem{ha2018}
D.~Ha and J.~Schmidhuber, ``World Models,'' \emph{arXiv:1803.10122}, 2018.

\bibitem{hafner2021}
D.~Hafner \emph{et~al.}, ``Mastering Atari with Discrete World Models,'' in \emph{ICLR}, 2021.

\bibitem{rai2020}
A.~Rai, ``Explainable AI: From Black Box to Glass Box,'' \emph{Journal of the Academy of Marketing Science}, vol.~48, pp.~137--141, 2020.

\bibitem{russell2019}
S.~Russell, \emph{Human Compatible: Artificial Intelligence and the Problem of Control}.\hskip 1em plus 0.5em minus 0.4em\relax Viking, 2019.

\bibitem{dignum2019}
V.~Dignum, \emph{Responsible Artificial Intelligence: How to Develop and Use AI in a Responsible Way}.\hskip 1em plus 0.5em minus 0.4em\relax Springer, 2019.

\bibitem{rahwan2019}
I.~Rahwan \emph{et~al.}, ``Machine Behaviour,'' \emph{Nature}, vol.~568, pp.~477--486, 2019.

\bibitem{yao2023}
S.~Yao \emph{et~al.}, ``ReAct: Synergizing Reasoning and Acting in Language Models,'' in \emph{ICLR}, 2023.

\bibitem{shinn2023}
N.~Shinn \emph{et~al.}, ``Reflexion: Language Agents with Verbal Reinforcement,'' \emph{arXiv:2303.11366}, 2023.

\bibitem{wang2023}
G.~Wang \emph{et~al.}, ``Voyager: An Open-Ended Embodied Agent with LLMs,'' \emph{arXiv:2305.16291}, 2023.

\bibitem{dignum2004}
F.~Dignum \emph{et~al.}, ``Meeting the Deadline: Why, When and How,'' in \emph{Formal Approaches to Agent-Based Systems}.\hskip 1em plus 0.5em minus 0.4em\relax Springer, 2004.

\bibitem{sergot2001}
M.~Sergot, ``A Computational Theory of Normative Positions,'' \emph{ACM Trans.\ on Computational Logic}, vol.~2, no.~4, pp.~581--622, 2001.

\bibitem{bordini2007}
R.~H.~Bordini, J.~F.~H{\"u}bner, and M.~Wooldridge, \emph{Programming Multi-Agent Systems in AgentSpeak using Jason}.\hskip 1em plus 0.5em minus 0.4em\relax Wiley, 2007.

\bibitem{simonsen2012}
J.~Simonsen and T.~Robertson, \emph{Routledge International Handbook of Participatory Design}.\hskip 1em plus 0.5em minus 0.4em\relax Routledge, 2012.

\bibitem{maler2004}
O.~Maler and D.~Nickovic, ``Monitoring Temporal Properties of Continuous Signals,'' in \emph{Formal Techniques, Modelling and Analysis of Timed and Fault-Tolerant Systems}.\hskip 1em plus 0.5em minus 0.4em\relax Springer, 2004.

\bibitem{platzer2018}
A.~Platzer, \emph{Logical Foundations of Cyber-Physical Systems}.\hskip 1em plus 0.5em minus 0.4em\relax Springer, 2018.

\bibitem{bengio2020}
Y.~Bengio \emph{et~al.}, ``A Meta-Transfer Objective for Learning to Disentangle Causal Mechanisms,'' in \emph{ICLR}, 2020.

\bibitem{abel2017}
D.~Abel \emph{et~al.}, ``Agent-Agnostic Human-in-the-Loop Reinforcement Learning,'' \emph{arXiv:1701.04079}, 2017.

\bibitem{finn2017}
C.~Finn \emph{et~al.}, ``Model-Agnostic Meta-Learning for Fast Adaptation,'' in \emph{ICML}, 2017.

\bibitem{euaiact2021}
European Commission, ``Proposal for a Regulation on Artificial Intelligence (AI Act),'' 2021.

\bibitem{bond2026ge}
[Anonymous], \emph{Geometric Ethics: The Mathematical Structure of Moral Reasoning}, v1.14, Department of Computer Engineering, [Institution], Spring 2026.

\bibitem{bond2026erisml}
[Anonymous], ``ErisML/DEME v3.0: A Library for Governed, Foundation-Model-Enabled Agents with Tensorial Ethics,'' [URL removed for blind review], 2026.

\bibitem{bond2026noescape}
[Anonymous], ``No Escape: Mathematical Containment for Artificial Agents---Why Structural Constraints Succeed Where Behavioral Rules Fail,'' [Institution], 2026.

\bibitem{bond2026noether}
[Anonymous], ``Noether's Theorem for Ethics: Harm Accounting and the Formal Structure of Normative Coherence,'' [Institution], 2026.

\bibitem{bond2026diffgeom}
[Anonymous], ``Differential Geometry for Moral Alignment: The Mathematical Foundations of DEME 3.0,'' [Institution], 2025.

\bibitem{bond2026stratified}
[Anonymous], ``Stratified Gauge Theory: Invariance and Conservation on Singular Spaces,'' [Institution], 2025.

\bibitem{bond2026deme2}
[Anonymous], ``DEME 2.0: Real-Time Ethical Governance for Safety-Critical Autonomous Systems via Computable Moral Landscapes,'' [Institution], 2025.

\bibitem{thiele2026}
[Anonymous collaborator], ``Does LaBSE Encode Moral Geometry? Testing the Geometric Ethics Framework Through Cross-Lingual Probing,'' Undergraduate Research Report, [Institution], 2026.

\end{thebibliography}

\appendix

\section{ErisML Grammar (Excerpt)}
\label{app:grammar}

\begin{lstlisting}
<model> ::= <environment> <agent>* <norms>?
            <dynamics>? <validation>?

<environment> ::= "environment" <id> "{" <env-body> "}"
<env-body> ::= <objects> <state> <observations>?
               <dynamics>

<agent> ::= "agent" <id> "{" <agent-body> "}"
<agent-body> ::= <capabilities> <beliefs> <intents>
                 <constraints>?

<norms> ::= "norms" <id> "{" <norm-rule>* "}"
<norm-rule> ::= <permission> | <prohibition>
              | <obligation> | <sanction>

<permission>  ::= "permission:" "{" <action-expr>
                    <condition>? "}" ";"
<prohibition> ::= "prohibition:" "{" <action-expr>
                    <condition>? "}" ";"
<obligation>  ::= "obligation:" "{" <action-expr>
                    <deadline>? "}" ";"

<intents> ::= "intents:" <intent-expr> ";"
<intent-expr> ::= <weighted> | <lexicographic>
                | <constrained>

<condition> ::= "when" <expr> | "unless" <expr>
\end{lstlisting}

\section{Compilation Example (PDDL)}
\label{app:pddl}

\textbf{ErisML Input}:

\begin{lstlisting}
environment SimpleRoom {
  objects: Robot, Door;
  state:
    door.status: {open, closed};
    robot.location: {inside, outside};
  dynamics:
    open_door() ~> door.status = open;
    enter() ~> robot.location = inside
               if door.status == open;
}

agent Robot {
  intents: goal robot.location == inside;
}

norms Safety {
  prohibition: Robot.enter()
    unless door.status == open;
}
\end{lstlisting}

\textbf{PDDL Output}:

\begin{lstlisting}
(define (domain simple-room)
  (:requirements :strips :typing)
  (:predicates (door-open) (door-closed)
               (robot-inside) (robot-outside))

  (:action open-door
    :parameters ()
    :precondition (door-closed)
    :effect (and (door-open)
                 (not (door-closed))))

  (:action enter
    :parameters ()
    :precondition (and (robot-outside)
                       (door-open)) ; Norm enforced
    :effect (and (robot-inside)
                 (not (robot-outside)))))

(define (problem reach-inside)
  (:domain simple-room)
  (:init (door-closed) (robot-outside))
  (:goal (robot-inside)))
\end{lstlisting}

\textbf{Note}: The prohibition from \texttt{norms Safety} is compiled
into the precondition of the \texttt{enter} action. This ensures the
planner never generates norm-violating plans, but provenance information
(that this precondition comes from a norm) is lost in translation.

\end{document}
